\usepackage[nopostdot, numberedsection, style=super, section, toc, acronyms, nogroupskip]{glossaries}
\usepackage{xparse}

\setlength{\glsdescwidth}{0.8\textwidth}
\renewcommand{\glsnamefont}[1]{\textbf{#1}}

% label, acronym, name, description
\DeclareDocumentCommand{\newdualentry}{ m m m m } {
	\newglossaryentry{gls-#1}{
	name={#3 (\glstext{#1})},
	description={#4}, nonumberlist
	}
	\newglossaryentry{#1}{
		type=\acronymtype,
		name={#2},
		first={#3 (#2)},
		firstplural={#3s (#2s)},
		description=\glslink{gls-#1}{#3},
		nonumberlist
	}%
}
%%% \gls{ABBREV} to point at Acronym/Abbreviation index
%%% \gls{gls-ABBREV} to point to glossary

\makeglossaries

\newdualentry{ADC}%
  {ADC}%
  {Analog to Digital Converter}%
  {An Analog to Digital Converter is a system that converts an analog signal
   into a quantized digital signal. Its counterpart is the \gls{DAC}.}%

\newdualentry{ABI}%
  {ABI}%
  {Application Binary Interface}%
  {An Application Binary Interface defines the low-level interface between
   an application and the operating system or between application components
   running on the same hardware.}%

\newdualentry{API}%
  {API}%
  {Application Programming Interface}%
  {The Application Programming Interface defines how a developer can write
   a program that requests services from an operating system or application.
   \glspl{API} are implemented by function calls composed of verbs and nouns,
   i.e. a function to execute on an object.}%

\newdualentry{ASan}%
  {ASan}%
  {AddressSanitizer}%
  {AddressSanitizer is a dynamic analysis tool that detects memory errors in %
   C and C++ programs, such as buffer overflows, use-after-free, and memory %
   leaks, by instrumenting the compiler.}%

\newdualentry{BSP}%
  {BSP}%
  {Board Support Package}%
  {A Board Support Package is the implementation of a specific interface defined
   by the abstract layer of an operating system that enables the latter to run
   on the particular hardware platform.}%

\newdualentry{CERT}%
  {CERT}%
  {Computer Emergency Response Team}%
  {The Computer Emergency Response Team is the coordination centre at the %
   Software Engineering Institute that investigates and reports on software %
   security issues.}%

\newdualentry{CTU}%
  {CTU}%
  {Cross Translation Unit}%
  {Cross Translation Unit analysis is an extension to static analysis that %
   follows calls and references across translation unit boundaries rather than %
   analysing a single translation unit in isolation.}%

\newdualentry{CVSS}%
  {CVSS}%
  {Common Vulnerability Scoring System}%
  {The Common Vulnerability Scoring System is a standard for rating the %
   severity of software vulnerabilities.}%

\newdualentry{CWE}%
  {CWE}%
  {Common Weakness Enumeration}%
  {The Common Weakness Enumeration is a community developed list of common %
   software and hardware weakness types.}%

\newdualentry{CPU}%
  {CPU}%
  {Central Processing Unit}%
  {The Central Processing Unit is the electronic circuitry that interprets
  instructions of a computer program and performs control logic, arithmetic,
  and input/output operations specified by the instructions. It maintains
  high-level control of peripheral components, such as memory and other devices.}%

\newdualentry{DAC}%
  {DAC}%
  {Digital to Analog Converter}%
  {A Digital to Analog Converter is a system that converts a quantized digital
   signal into an analog signal. Its counterpart is the \gls{ADC}.}%

\newdualentry{DMA}%
  {DMA}%
  {Direct Memory Access}%
  {Direct Memory Access is a feature of a computer system that allows hardware
   subsystems to access main system \gls{gls-RAM} directly, thereby bypassing
   the \gls{gls-CPU}.}%

\newacronym{DRD}{DRD}{Document Requirements Definition}

\newacronym{CIDL}{CIDL}{Configuration Item Data List}

\newacronym{ADD}{ADD}{Architectural Design Document}

\newacronym{CCB}{CCB}{Configuration Control Board}

\newacronym{IRD}{IRD}{Interface Requirements Document}

\newacronym{PQA}{PQA}{Product Quality Assurance Plan}

\newacronym{QTS}{QTS}{Qualification Test Specification}

\newdualentry{DSP}%
  {DSP}%
  {Digital Signal Processor}%
  {A Digital Signal Processor is a specialised processor with its architecture %
   targeting the operational needs of digital signal processing.}%

\newglossaryentry{Doxygen}{
  name={Doxygen},
  description={Doxygen is a documentation generator for writing software %
   	       reference documentation.},
  nonumberlist
  }

\newdualentry{ELF}%
  {ELF}%
  {Executable and Linkable Format}%
  {The Executable and Linkable Format is a common standard file format for
   executables, object code, shared libraries, and core dumps.}%

\newacronym{FDIR}{FDIR}{Fault Detection, Isolation and Recovery}

\newdualentry{EDAC}%
  {EDAC}%
  {Error Detection And Correction}%
  {Error Detection And Correction is a technique used to detect and correct
   errors in digital data, typically in memory, through the use of redundant
   check bits. On fault-tolerant SPARC processors it is used to protect memory
   against single-event upsets.}%

\newdualentry{EDF}%
  {EDF}%
  {Earliest Deadline First}%
  {Earliest Deadline First is a dynamic priority scheduling algorithm that
   assigns priorities based on task deadlines. Tasks with the nearest
   deadline receive the highest priority.}%

\newdualentry{FIFO}%
  {FIFO}%
  {First In - First Out}%
  {In FIFO processing, the "head" element of a queue is processed first.
   Once complete, the element is removed and the next element in line becomes
   the new queue head.}%

\newglossaryentry{FlightOS}{
  name={FlightOS},
  description={FlightOS is a lightweight real-time operating system for
               space-based astronomy applications, targeting European
               \gls{SPARC}-based processors. It is licensed under the
               \gls{GPL}.},
  nonumberlist
  }

\newdualentry{FPU}%
  {FPU}%
  {Floating Point Unit}%
  {A co-processor unit that specialises in floating-point calculations.}%

\newdualentry{ILP}%
  {ILP}%
  {Instruction Level Parallelism}%
  {Instruction-level parallelism (ILP) is a measure of how many instructions in
   a computer program can be executed simultaneously by the \gls{CPU}.}%

\newdualentry{ISR}%
  {ISR}%
  {Interrupt Service Routine}%
  {An Interrupt Service Routine is a function that handles the actions needed
   to service an interrupt.}%

\newacronym{FSW}{FSW}{Flight Software}

\newdualentry{GCC}%
  {GCC}%
  {GNU Compiler Collection}%
  {The GNU Compiler Collection is a compiler system produced by the
   GNU project. It is part of the GNU toolchain collection of programming
   tools.}%

\newdualentry{GPL}%
  {GPL}%
  {GNU General Public License}%
  {The GNU General Public License is a free software licence. Version 2 %
   (GPLv2) is used to license FlightOS and, unless stated otherwise, the %
   documentation examples reproduced in this document.}%

\newglossaryentry{GNU}{
  name={GNU},
  description={GNU is a recursive acronym that stands for "GNU's Not Unix!".
	      The GNU project is a free software collaboration project announced
	      in 1987. Users are free to run GNU-licensed software, share, copy,
	      distribute, study and modify it. GNU software guarantees these
	      freedom-rights legally via its license and is thefore free
      	      software.},
  nonumberlist
}

\newglossaryentry{LEON2}{
  name={LEON2},
  description={The LEON2 is a synthesisable VHDL model of a 32-bit processor
	       compliant with the SPARC V8 architecture. It is highly
	       configurable and particularly suitable for \gls{SoC} designs.
	       Its source code is available under the GNU LGPL license},
  nonumberlist
  }

\newglossaryentry{LEON3}{
  name={LEON3},
  description={The LEON3 is an updated version of the \gls{LEON2}, changes
	       include \gls{gls-SMP} support and a deeper instruction pipeline},
  nonumberlist
  }

\newglossaryentry{LEON3-FT}{
  name={LEON3-FT},
  description={The LEON3-FT is a fault-tolerant version of the \gls{LEON3}.
  	       Changes to the base version include autonomous error handling,
       	       cache locking and different cache replacement strategies.},
  nonumberlist
  }

\newglossaryentry{LEON4}{
  name={LEON4},
  description={The LEON4 is an updated version of the \gls{LEON3} family of
               SPARC V8 processors, offering a deeper pipeline, higher clock
               frequencies and a wider memory interface.},
  nonumberlist
  }

\newglossaryentry{GR740}{
  name={GR740},
  description={The GR740 is a system-on-chip based on four \gls{LEON4FT}
               cores with \gls{SMP} support, manufactured by Gaisler
               \cite{GR740UM}.},
  nonumberlist
  }

\newglossaryentry{GR712RC}{
  name={GR712RC},
  description={The GR712RC is a dual-core \gls{LEON3-FT} system-on-chip
               manufactured by Gaisler \cite{GR712UM}.},
  nonumberlist
  }

\newglossaryentry{LEON4FT}{
  name={LEON4FT},
  description={The LEON4FT is the fault-tolerant version of the \gls{LEON4}
               processor core, adding error detection and correction
               capabilities.},
  nonumberlist
  }

\newglossaryentry{ramses}{
  name={RAMSES},
  description={The RAMSES project is a space mission whose on-board data
               processing platform is based on a custom built computing device
               using the \gls{GR712RC} processor.},
  nonumberlist
  }

\newglossaryentry{ariel}{
  name={ARIEL},
  description={The ARIEL (Atmospheric Remote-sensing Infrared Exoplanet
               Large-survey) project is a space mission whose on-board data
               processing platform is based on a custom built computing device
               using the \gls{GR712RC} processor.},
  nonumberlist
  }

\newglossaryentry{smile}{
  name={SMILE},
  description={The SMILE (Solar wind Magnetosphere Ionosphere Link Explorer)
               project is a space mission whose on-board data processing
               platform is based on a custom built computing device using the
               \gls{GR712RC} processor.},
  nonumberlist
  }

\newdualentry{MMU}%
  {MMU}%
  {Memory Management Unit}%
  {A Memory Management Unit performs address space translation between physical
   and virtual memory pages and protects unprivileged access to certain memory
   regions.}%

\newdualentry{MPPB}%
  {MPPB}%
  {Massively Parallel Processor Breadboarding system}%
  {The Massively Parallel Processor Breadboarding system is a proof-of-concept %
   design for a space-hardened, fault-tolerant multi-DSP system with various %
   subsystems to build a powerful digital signal processing system with a high %
   data throughput. Its distinguishing features are the \gls{gls-NoC} and the
   \gls{Xentium} \glspl{DSP} controlled by a \gls{LEON2} processor.
   It was developed under ESA contract 21986 by Recore Systems B.V.}%

\newacronym{NCR}{NCR}{Non-Conformance Reports}

\newdualentry{NGAPP}%
  {NGAPP}%
  {Next Generation Astronomy Processing Platform}%
  {Next Generation Astronomy Processing Platform was an evaluation of the
   \gls{MPPB} performed in a joint effort of RUAG Space Austria and the
   Department of Astrophysics of the University of Vienna.
   The project was funded under ESA contract 40000107815/13/NL/EL/f.}%

\newdualentry{NoC}%
  {NoC}%
  {Network On Chip}%
  {A Network On Chip is a communication system on an integrated circuit that
   applies (packet based) networking to on-chip communication. It offers
   improvements over more conventional bus interconnects and is more scalable
   and power efficient in complex \gls{gls-SoC} desgins.}%

\newacronym{NRB}{NRB}{Nonconformance Review Board}

\newacronym{PA}{PA}{Product Assurance}

\newacronym{PAM}{PAM}{Product Assurance Manager}

\newdualentry{POSIX}
  {POSIX}
  {Portable Operating System Interface}
  {The Portable Operating System Interface is a family of standards specified
   by the IEEE Computer Society for maintaining compatibility between
   operating systems.}%

\newdualentry{PUS}
  {PUS}
  {Packet Utilisation Standard}
  {The Packet Utilisation Standard addresses the end-to-end transport of telemetry
   and telecommand data between user applications on the ground and applications
   onboard a satellite. See also ECSS-E-70-41A.}%

\newdualentry{RAM}%
  {RAM}%
  {Random-Access Memory}%
  {Random-Access Memory is a type of memory where each memory cell may be
   accessed directly via their memory addresses.}%

\newacronym{RID}{RID}{Review Item Discrepancy}

\newdualentry{RISC}%
  {RISC}%
  {Reduced Instruction Set Computing}%
  {RISC is a \gls{CPU} design strategy that intends to improve performance by
   combining a simplified instruction set with a microprocessor architecture
   that is capable of executing an instruction in a smaller number of clock
   cycles.}%

\newdualentry{RMAP}%
  {RMAP}%
  {Remote Memory Access Protocol}%
  {The Remote Memory Access Protocol is a form of \gls{SpaceWire} communication
   that transparently communicates writes to memory mapped regions between
   different hardware devices.}%


\newdualentry{ROSE}%
  {ROSE}%
  {ROSE compiler framework}%
  {The ROSE compiler framework is an open source compiler infrastructure that %
   provides a rich abstract syntax tree representation of source code, on which %
   static analysis tools such as Rosecheckers are built.}%

\newdualentry{RR}%
  {RR}%
  {Round Robin}%
  {Round Robin is a scheduling algorithm where time slices are assigned in equal
   poritions and in circular order. In the context of threads, priorities are
   usually only used to control re-scheduling order when a mutex is accessed by
   a thread.}%

\newdualentry{RTOS}%
  {RTOS}%
  {Real-Time Operating System}%
  {A Real-Time Operating System is an operating system intended to serve
   real-time applications that process data as it comes in, typically
   without buffer delays, and guarantees a certain capability within a
   specified time constraint.}%

\newacronym{RSA}{RSA}{RUAG Space Austria}

\newacronym{SCF}{SCF}{Software Configuration File}

\newacronym{SDD}{SDD}{Software Design Document}

\newacronym{SDP}{SDP}{Software Development Plan}

\newacronym{SMCP}{SMCP}{Software Management and Configuration Plan}

\newacronym{SPSC}{SPSC}{Single-Producer Single-Consumer}

\newdualentry{SMP}%
  {SMP}%
  {Symmetric Multiprocessing}%
  {Symmetric Multiprocessing denotes computer architectures, where two or more
   identical processors are connected to the same periphery and are controlled
   by the same operating system instance.}%

\newdualentry{SEI}%
  {SEI}%
  {Software Engineering Institute}%
  {The Software Engineering Institute is a research and development centre at %
   Carnegie Mellon University that advances software engineering practice and %
   publishes the SEI CERT Coding Standards.}%

\newacronym{SPAMR}{SPAMR}{Software Product Assurance Milestone Report}

\newacronym{SPR}{SPR}{Software Problems Reports}

\newacronym{SSS}{SSS}{System Software Specification}

\newacronym{SRD}{SRD}{Software Release Document}

\newacronym{SRS}{SRS}{Software Requirement Specification}

\newdualentry{SoC}%
  {SoC}%
  {System On Chip}%
  {A System On Chip is an integrated circuit that combines all components of a %
   computer or other electronic system into a single chip.}%

\newdualentry{SRMMU}%
  {SRMMU}%
  {SPARC Reference Memory Management Unit}%
  {The SPARC Reference Memory Management Unit is the reference implementation
   of an MMU for the \gls{SPARC} architecture. It provides virtual memory
   translation and memory protection.}%

\newglossaryentry{SpaceWire}{
  name={SpaceWire},
  description={SpaceWire is a spacecraft communication network based in part
               on the IEEE 1355 standard of communications.},
  nonumberlist
  }

\newglossaryentry{SPARC}{
  name={SPARC},
  description={SPARC ("scalable processor architecture") is a \gls{gls-RISC}
  	       instruction set architecture developed by Sun Microsystems in
	       the 1980s. The distinct feature of SPARC processors is the high
	       number of \gls{gls-CPU} registers that are accessed similarly to
	       stack variables via ``sliding windows''.},
  nonumberlist
  }


\newacronym{SQ}{SQ}{Software Quality}

\newacronym{SQAM}{SQAM}{Software Quality Assurance Manager}

\newdualentry{SSDP} % label
  {SSDP}            % abbreviation
  {Scalable Sensor Data Processor}  % long form
  {The Scalable Sensor Data Processor (SSDP) was a next generation on-board %
   data processing mixed-signal ASIC, envisaged for use in scientific payloads %
   requiring high-performance on-board processing capabilities. It was built on %
   a heterogeneous multicore architecture, combining two %
   \gls{Xentium} \gls{DSP} cores with a general-purpose \gls{LEON3-FT} control %
   processor in a \gls{gls-NoC}.} % description

\newacronym{SQA}{SQA}{Software Quality Assurance}

\newdualentry{TCM}%
  {TCM}%
  {Tightly-Coupled Memory}%
  {Tightly-Coupled Memory is the local data memory that is directly accessible %
   by a Xentium's load/store unit. It can be viewed as a completely %
   program-controlled data cache.}%

\newdualentry{TOCTOU}%
  {TOCTOU}%
  {Time Of Check To Time Of Use}%
  {Time Of Check To Time Of Use is a vulnerability where a resource is checked
   for a condition at one point in time and then used at a later point in time,
   when the condition may have changed.}%


\newacronym{TP}{TP}{Test Plan}
\newacronym{TS}{TS}{Test Specification}

\newacronym{SVR}{SVR}{Software Verification Report}

\newacronym{VQR}{VQR}{Software Verification and Quality Report}

\newdualentry{UBSan}%
  {UBSan}%
  {UndefinedBehaviorSanitizer}%
  {UndefinedBehaviorSanitizer is a fast undefined behavior detector that %
   modifies the program at compile-time to catch various kinds of undefined %
   behavior during program execution.}%

\newdualentry{UVIE}%
  {UVIE}%
  {University of Vienna}%
  {University of Vienna.}%

\newdualentry{VLIW}%
  {VLIW}%
  {Very Long Instruction Word}%
  {Very Long Instruction Word is a processor architecture design concept that
   exploits \gls{gls-ILP}. This approach allows higher performance at a smaller
   silicone footprint compared to serialised instruction processors, as no
   instruction re-ordering logic to exploit superscalar capabilities of the
   processor must be integrated on the chip, but requires either code to be
   tuned manually or a very sophisticated compiler to exploit the full potential
   of the processor.}%

\newdualentry{WCET}%
  {WCET}%
  {Worst-Case Execution Time}%
  {Worst-Case Execution Time is the maximum length of time that a task or
   function can take to execute on a specific hardware platform. It is a
   critical parameter in the design and analysis of real-time systems.}%

\newglossaryentry{Xentium}{
  name={Xentium},
  description={The Xentium is a high performance \gls{gls-VLIW} \gls{DSP} core.
               It operates 10 parallel execution slots supporting 32/40 bit
	       scalar and two 16-bit element vector operations.},
  nonumberlist
  }


\newdualentry{CCSDS}%
  {CCSDS}%
  {Consultative Committee for Space Data Systems}%
  {The Consultative Committee for Space Data Systems is the body that publishes %
   the Space Packet Protocol used as the base transport of \gls{PUS} packets.}%

\newdualentry{CRC}%
  {CRC}%
  {Cyclic Redundancy Check}%
  {A Cyclic Redundancy Check is an error-detecting code commonly used to detect %
   accidental changes to digital data, e.g. in the \gls{RMAP} and \gls{PUS} %
   protocols.}%

\newdualentry{EGSE}%
  {EGSE}%
  {Electrical Ground Support Equipment}%
  {Electrical Ground Support Equipment is the collection of test and support %
   equipment used to verify a space segment on ground, e.g. the CHEOPS model %
   operated through the SpaceWire bridge.}%

\newdualentry{EOP}%
  {EOP}%
  {End Of Packet}%
  {End Of Packet is the \gls{SpaceWire} marker that terminates a packet, the %
   counterpart of the Error End Of Packet (EEP) marker that signals an %
   aborted transmission.}%

\newdualentry{FEE}%
  {FEE}%
  {Front-End Electronics}%
  {Front-End Electronics is the electronics directly interfacing detectors and %
   processing their signals; on CHEOPS the FEE produces the housekeeping and %
   data frames that the bridge can interpret.}%

\newglossaryentry{GRESB}{
  name={GRESB},
  description={The GRESB is a SpaceWire/Ethernet bridge by Gaisler that routes %
	       \gls{SpaceWire} packets to and from TCP sockets over an Ethernet %
	       network; it offers three SpaceWire links and six virtual links %
	       \cite{GRESBUM}.},
  nonumberlist
  }

\newglossaryentry{STAR}{
  name={STAR API},
  description={The STAR API is the software interface to SpaceWire interface %
	       devices such as the Brick MK II, Brick MK III, PCI and PCIe %
	       cards; \texttt{spw\_bridge} uses it for all device, channel and %
	       link management.},
  nonumberlist
  }

\newglossaryentry{SpW}{
  name={SpW},
  description={SpW is a common abbreviation of \gls{SpaceWire} used throughout %
	       this document.},
  nonumberlist
  }

\newdualentry{TCP}%
  {TCP}%
  {Transmission Control Protocol}%
  {Transmission Control Protocol is the stream-oriented transport protocol of %
   the Internet protocol suite; the bridge carries \gls{SpaceWire} traffic %
   over a single TCP connection by default.}%

\newdualentry{UDP}%
  {UDP}%
  {User Datagram Protocol}%
  {User Datagram Protocol is the datagram-oriented transport protocol of the %
   Internet protocol suite; with \texttt{-u} the bridge maps a single UDP %
   datagram to a single \gls{SpaceWire} packet.}%

\newdualentry{ECSS}%
  {ECSS}%
  {European Cooperation for Space Standardization}%
  {The European Cooperation for Space Standardization issues the standards %
   series used by the European space industry, e.g. ECSS-E-ST-70-41C for the %
   \gls{PUS} and ECSS-E-ST-50-52C for \gls{RMAP}.}%

\newglossaryentry{MSSL}{
  name={MSSL},
  description={The Mullard Space Science Laboratory is a department of %
	       University College London; on CHEOPS it provides the science %
	       payload, whose FEE data frames follow MSSL-IF-109 %
	       \cite{MSSLIF109}.},
  nonumberlist
  }


% do not remove
\glsresetall
\makeglossaries
